\scanpage{15}
% Foliotation archivistique au crayon, en bas à gauche : 14.
\deleted{$k=\text{corps premier}$}

\deleted{$K$ \unsure{corps}\quad $k=\text{corps premier}$}, $k=\mathbb Q$

$\widehat{\mathcal X}$ droite projective sur $\bar K$
\[
S\subset\mathcal X(K)\qquad S\text{ partie finie}
\]
\[
\mathcal X=\widehat{\mathcal X}\setminus S
\]
$\mathcal X'$ rev. abélien [modéré] maximal de $\mathcal X$.

A) Si $\operatorname{Card}S\in\{0,1\}$, alors $\mathcal X'=\mathcal X$ \ldots

B) Si $\operatorname{Card}S=2$, alors $\mathcal X$ est \inserted{(canonique)} un torseur sous une forme de $\mathbb G_{m\,\overline{\mathbb Q}}$;
\deleted{et} choisir un des pts de $S$ comme $\infty$ \inserted{(l'autre étant pris comme $0$)} revient à choisir un isom du \deleted{\unsure{ce}} groupe d'op. avec $\mathbb G_m$; choisir une ``\uline{origine}'' (sur $1$) de $\mathcal X$ revient à trivialiser ce torseur.
\deleted{Ceci étant faits, on trouve} \inserted{\unsure{En tt cas}} on a une descente de la situation $(\widehat{\mathcal X},\mathcal X)$ de $\overline{\mathbb Q}$ à $\mathbb Q$ en $(\widehat X,X)$
\[
\widehat X=\mathbb P^1_{\mathbb Q},\quad X=\mathbb P^1_{\mathbb Q}\setminus\{0,\infty\}=\mathbb G_{m\,\mathbb Q}.
\]
Ici
\[
\mathcal G=\Gamma\mathbin{\underset{\scriptscriptstyle 1/2}{.}}\overline{\mathbb Q}^{*}
\]
$\Gamma$ opérant sur $\overline{\mathbb Q}^{*}$ de la façon évidente. Le rev. kummérien\inserted{-type} d'indice $n$ de $X$ est $X_n\inserted{=X}$, avec
\diagram{\begin{tikzcd}[ampersand replacement=\&]
X_n \arrow[r,"{t\mapsto t^n}"] \& X
\end{tikzcd}},
rev. principal de groupe $\mu_{n\,\mathbb Q}=\mu_n$.
\margin{\diagram{\begin{tikzcd}[ampersand replacement=\&]
X_n \arrow[d,"{f_n(t)=t^n}"] \\ X
\end{tikzcd}}}
On a, pour
\diagram{\begin{tikzcd}[ampersand replacement=\&,row sep=small]
\mathcal G_n=\operatorname{Aut}_k(\mathcal X_n\to\mathcal X) \arrow[d,"\simeq"'] \\
\operatorname{Aut}_k(\mathcal X_n)
\end{tikzcd}}
\[
\mathcal G'_n=\Gamma\mathbin{\underset{\scriptscriptstyle 1/2}{.}}\overline{\mathbb Q}^{*}.
\]
L'homomorphisme canon.
\diagram{\begin{tikzcd}[ampersand replacement=\&]
\mathcal G_n \arrow[r] \& \mathcal G
\end{tikzcd}}
donne l'identité sur $\deleted{\unsure{\mathcal G}}\ \Gamma$, et l'élévation à la puiss. nième sur $\unsure{\mathbb G|G}$.
